\chapter{Chain Conditions}

Let $\Sigma$ be a partially ordered set. $\Sigma$ said to be satisfying the \textcolor{orange}{ascending chain condition} (\textcolor{orange}{a.c.c}) if, for every increasing sequence $x_1\leqslant x_2\leqslant\cdots$ in $\Sigma$ there exists $n$ such that $x_n=x_{n+1}=\cdots$.

Similarly, $\Sigma$ said to be satisfying the \textcolor{orange}{descending chain condition} (\textcolor{orange}{d.c.c}) if, for every decreasing sequence $x_1\geqslant x_2\geqslant\cdots$ in $\Sigma$ there exists $n$ such that $x_n=x_{n+1}=\cdots$.

A module satisfying the ascending chain condition is said to be \textcolor{orange}{Noetherian}. A module satisfying the descending chain condition is said to be \textcolor{orange}{Artinian}.

\begin{proposition}{6.2}
	$M$ is a Noetherian $A$-module $\Leftrightarrow$ every submodule of $M$ is finitely generated.
\end{proposition}

\begin{proposition}{6.3}
	Let $0\rightarrow M^\prime\rightarrow M\rightarrow M^{\prime\prime}\rightarrow0$ be an exact sequence of $A$-modules. Then
	\begin{enumerate}[label=\textup{(\roman*)}]
		\item $M$ is Noetherian $\Leftrightarrow$ $M^\prime$ and $M^{\prime\prime}$ are Noetherian.
		\item $M$ is Artinian $\Leftrightarrow$ $M^\prime$ and $M^{\prime\prime}$ are Artinian.
	\end{enumerate}
\end{proposition}

\begin{proposition}{6.5}
	Let $A$ be a Noetherian (resp. Artinian) ring, $M$ a finitely generated $A$-module. Then $M$ is Noetherian (resp. Artinian).
\end{proposition}

\begin{definition}
	A \textcolor{orange}{chain} of submodules of a module $M$ is a sequence $\left(M_i\right)$ of submodules of $M$ such that
	\begin{equation*}
		M=M_0\supsetneq M_1\supsetneq\cdots\supsetneq M_n=0.
	\end{equation*}
	The \textcolor{orange}{length} of the chain is $n$. A \textcolor{orange}{composition series} of $M$ is a maximal chain, that is one which no extra submodules can be inserted, or equivalently $M_{i-1}/M_i$ is simple.
\end{definition}

\begin{proposition}{6.7}
	Suppose that $M$ has a composition series of length $n$. Then every composition series of $M$ has length $n$, and every chain can be extended to a composition series.
\end{proposition}

The length of a composition series of $M$ is called the \textcolor{orange}{length} of $M$, denoted as $l(M)$.

\begin{proposition}{6.10}
	For a vector space, 
	\begin{center}
		finite dimension $\Longleftrightarrow$ finite length $\Longleftrightarrow$ a.c.c. $\Longleftrightarrow$ d.c.c..
	\end{center}
	If the conditions are satisfied, length equals dimension.
\end{proposition}

\section*{DEFINITIONS AND PROPOSITIONS FROM EXERCISES}

A topological space $X$ is said to be \textcolor{orange}{Noetherian} if the open subsets of $X$ satisfy the ascending chain condition. Or equivalently, the closed subsets satisfy descending chain condition.